\section{课程概览与生成模型的核心问题}

本讲是 KAIST CS492D ``Diffusion and Flow Models'' 课程的第一讲。在正式进入 Diffusion Model 和 Flow Model 之前，授课教师 Minhyuk Sung 首先回顾了生成模型（Generative Models）的基础知识，涵盖采样理论、GAN（Generative Adversarial Network）和 VAE（Variational Autoencoder）两大经典范式，以及支撑这些模型的概率统计工具。

\subsection{什么是生成模型？}

生成模型的核心任务可以用一句话概括：\textbf{给定一组数据样本，学习一个模型，使之能够生成新的、看起来与训练数据相似但又不完全相同的样本}。

以图像为例，假设我们有一个包含 58 亿张图像的超大规模数据集（如 LAION-5B），生成模型的目标就是训练一个模型，使其能够产出全新的、逼真的图像——既不是简单地从数据集中复制某张图片，也不是随机噪声，而是分布上与真实图像一致的全新样本。

\begin{importantbox}{生成模型的统计视角}
从统计学角度看，生成模型的问题可以形式化为：假设所有数据点 $x_1, x_2, \ldots, x_N$ 都是从某个未知的概率分布 $p_{\text{data}}(x)$ 中独立同分布（i.i.d.）采样得到的，我们的目标是学习这个分布，从而能够从 $p_{\text{data}}(x)$ 中采样新的数据点。
\end{importantbox}

为了帮助理解，授课教师从一个更简单的场景出发：假设我们有一组二维平面上的点，它们形成某种特定的形状（例如两个半月形）。那么问题变为：如何构建一个模型，使其能够生成新的点，使这些新点看起来像是从同一个分布中采样得到的？

\subsection{从简单到复杂：图像作为高维向量}

一张 $256 \times 256$ 的 RGB 图像，可以展平为一个 $256 \times 256 \times 3 = 196{,}608$ 维的向量。因此，一组图像就是高维空间中的一组点，图像生成的问题本质上等价于在这个高维空间中采样。

\begin{knowledgebox}{图像的向量化表示}
每张 $H \times W$ 的 RGB 图像可以看作 $\mathbb{R}^{H \times W \times 3}$ 空间中的一个点。图像集合就是该高维空间中的点云，它们共同定义了一个数据分布 $p_{\text{data}}(x)$。生成模型的任务就是从这个分布中采样新的点（即新的图像）。
\end{knowledgebox}

\subsection{本章小结}

生成模型的核心思想是将数据建模为某个概率分布的样本，通过学习这个分布来生成新数据。无论是二维点还是高维图像，问题的数学本质是一样的——从未知分布中采样。

\newpage
